% ============================================================================
% Fichier  : partie2/P02_C05_siem.tex
% Titre    : Le SIEM
% Auteur   : MINKA MI NGUIDJOÏ Thierry Emmanuel
% ============================================================================
% Sources : NIST SP 800-92 ; CIS Controls v8 ; MITRE ATT&CK.
% ============================================================================

\chapter{Le SIEM}
\label{chap:siem}
\chaptermark{Le SIEM}

\epigraph{%
    « Le SIEM ne détecte pas~; il rapproche. Ce sont vos règles qui détectent. »%
}{Principe du cours}

\niveauChapitre{Intermédiaire}{Chapitre~\ref{chap:journaux}.}

\begin{boxobjectifs}
\begin{itemize}
    \item Définir le SIEM et son rôle dans le SOC.
    \item Décrire la chaîne collecte, normalisation, corrélation, alerte.
    \item Comprendre la notion de cas d'usage de détection.
    \item Distinguer corrélation et simple recherche.
    \item Éviter les pièges du volume et du bruit.
\end{itemize}
\end{boxobjectifs}

\bigskip

% ============================================================================
\section*{Question directrice : un SIEM aurait-il, à lui seul, arrêté l'attaque~?}
% ============================================================================

Beaucoup croient qu'installer un SIEM « rend visible » l'attaque. En réalité,
un SIEM sans cas d'usage est un entrepôt de journaux muet. Ce chapitre montre
ce qu'un SIEM fait, et ce qu'il ne fait pas sans travail humain.

\medskip
\textbf{Plan du chapitre.}
\begin{enumerate}
    \item Qu'est-ce qu'un SIEM (section~\ref{sec:c05-definir}).
    \item La chaîne de traitement (section~\ref{sec:c05-chaine}).
    \item Les cas d'usage (section~\ref{sec:c05-cas}).
    \item Corréler, pas seulement chercher (section~\ref{sec:c05-correler}).
    \item Les pièges (section~\ref{sec:c05-pieges}).
\end{enumerate}

\separateur

% ============================================================================
\section{Qu'est-ce qu'un SIEM}
\label{sec:c05-definir}
% ============================================================================

\begin{boxdef}
\termecle{SIEM} (\textit{Security Information and Event Management})~:
plateforme qui collecte les journaux de sources multiples, les normalise, les
conserve, et déclenche des alertes lorsque des règles de corrélation
reconnaissent un motif suspect.
\end{boxdef}

% ============================================================================
\section{La Chaîne de Traitement}
\label{sec:c05-chaine}
% ============================================================================

\begin{tableaucours}{L{3.0cm}L{9.2cm}}{%
    \tableauHeader{Étape} &
    \tableauHeader{Rôle}
}
Collecte & Recevoir les journaux des sources (agents, syslog, API) \\
Normalisation & Ramener à un format commun (qui, quoi, quand, où) \\
Enrichissement & Ajouter contexte (géolocalisation, criticité de l'actif) \\
Corrélation & Rapprocher plusieurs événements selon une règle \\
Alerte & Signaler à l'analyste un motif suspect, avec son contexte \\
\end{tableaucours}
\vspace{-0.8em}
\begin{center}{\small\itshape Tableau~5.1. La chaîne de traitement d'un SIEM}\end{center}

% ============================================================================
\section{Les Cas d'Usage}
\label{sec:c05-cas}
% ============================================================================

\begin{boxdef}
\termecle{Cas d'usage de détection}~: situation précise que l'on veut
détecter, traduite en une ou plusieurs règles sur des sources identifiées.
Exemple~: « un compte de service se connecte à un poste de travail hors
plage horaire habituelle ».
\end{boxdef}

Un SIEM se construit cas d'usage par cas d'usage, en priorité sur les
techniques d'attaque les plus probables (ATT\&CK). Sans cas d'usage, il ne
produit aucune alerte utile.

% ============================================================================
\section{Corréler, pas seulement Chercher}
\label{sec:c05-correler}
% ============================================================================

Chercher dans les journaux (« trouve-moi les connexions de ce compte »)
répond à une question que l'on se pose déjà. \termecle{Corréler}, c'est
rapprocher automatiquement des événements de sources différentes pour faire
émerger une alerte que personne n'a pensé à chercher~: une connexion, suivie
d'une création de compte, suivie d'un accès à un partage sensible, en
quelques minutes.

% ============================================================================
\section{Les Pièges}
\label{sec:c05-pieges}
% ============================================================================

\begin{boxalert}
\textbf{Erreur fréquente.} Tout collecter et activer toutes les règles par
défaut.

\smallskip
\textbf{Règle~:} le volume coûte cher et le bruit noie les vraies alertes. On
collecte ce qui sert des cas d'usage, on règle les seuils, on mesure le taux
de faux positifs (Chapitre~\ref{chap:gouvernance}).
\end{boxalert}

\begin{boxscenario}{Opération BÉNOUÉ. Le SIEM arrivé trop tard}
\textbf{Constatation.} Après l'incident, LiZbiZ déploie un SIEM mutualisé
chez NYONG SÉCURITÉ. Rejouées sur les journaux de février, trois règles
simples (connexion anormale d'un compte de service, création de compte hors
maintenance, service installé à distance) auraient alerté dès le 23~février.

\medskip
\textbf{Analyse} (confiance élevée). Le SIEM n'aurait rien fait seul~: ce sont
ces trois cas d'usage, écrits à la main, qui auraient réduit le délai de
présence de onze jours à quelques heures.
\end{boxscenario}

\begin{boxmetier}{Analyste SOC}
\textbf{Un SIEM se nourrit de cas d'usage.} L'analyste ne « regarde » pas le
SIEM~: il y écrit des cas d'usage, les éprouve, mesure leur bruit et les
affine. La valeur est dans les règles, pas dans l'outil.
\end{boxmetier}

\separateur

\begin{boxresume}
\textbf{Ce qu'il faut retenir de ce chapitre~:}
\begin{itemize}
    \item \textbf{Le SIEM} collecte, normalise, corrèle, alerte.
    \item \textbf{Sans cas d'usage}, il est muet.
    \item \textbf{Corréler} fait émerger ce qu'on n'aurait pas cherché.
    \item \textbf{Le volume et le bruit} sont les pièges.
    \item \textbf{La valeur} est dans les règles, pas dans l'outil.
\end{itemize}
\end{boxresume}

\begin{boxinfo}
\textbf{Points de contrôle de l'élève.} Je sais~: définir un SIEM~; décrire
sa chaîne~; écrire un cas d'usage~; distinguer corréler et chercher~; éviter
le piège du volume.
\end{boxinfo}

\bigskip

\begin{boxexo}[title={\ding{45}\quad Exercice~5.1. Chaîne \niveaubase}]
Remettez dans l'ordre~: alerte, collecte, corrélation, normalisation.
\end{boxexo}

\begin{boxexo}[title={\ding{45}\quad Exercice~5.2. Muet \niveaubase}]
Pourquoi un SIEM sans cas d'usage ne produit-il aucune alerte utile~?
\end{boxexo}

\begin{boxexo}[title={\ding{45}\quad Exercice~5.3. Cas d'usage \niveaubase}]
Rédigez un cas d'usage en une phrase pour « compte de service sur un poste
ordinaire ».
\end{boxexo}

\begin{boxexo}[title={\ding{45}\quad Exercice~5.4. Corréler \niveaumoyen}]
Donnez un exemple de corrélation à trois événements qui révèle une intrusion.
\end{boxexo}

\begin{boxexo}[title={\ding{45}\quad Exercice~5.5. Bruit \niveaumoyen}]
Pourquoi activer toutes les règles par défaut nuit-il à la détection~?
\end{boxexo}

\begin{boxexo}[title={\ding{45}\quad Exercice~5.6. BÉNOUÉ \niveaumoyen}]
Quelles trois règles auraient alerté dès le 23~février~?
\end{boxexo}

\begin{boxexo}[title={\ding{45}\quad Exercice~5.7. Enrichissement \niveaumoyen}]
Qu'apporte l'enrichissement (criticité de l'actif) à une alerte~?
\end{boxexo}

\begin{boxexo}[title={\ding{45}\quad Exercice~5.8. Chercher ou corréler \niveaumoyen}]
Classez~: « liste les connexions de svc\_sauvegarde »~; « alerte si un compte
de service accède à un partage sensible après une connexion anormale ».
\end{boxexo}

\begin{boxexo}[title={\ding{45}\quad Exercice~5.9. Catalogue \niveauexpert}]
Proposez un catalogue de dix cas d'usage prioritaires pour LiZbiZ, reliés à
des techniques ATT\&CK.
\end{boxexo}

\begin{boxexo}[title={\ding{45}\quad Exercice~5.10. Dimensionnement \niveauexpert}]
Comment arbitrer entre volume collecté, coût et couverture de détection~?
\end{boxexo}

\bigskip

\begin{boxinfo}
\textbf{Pour aller plus loin.}
\begin{itemize}
    \item CIS Controls v8~\cite{cis2021controls}~; MITRE ATT\&CK~\cite{mitreattack}.
\end{itemize}
\end{boxinfo}

\begin{boxinfo}
\textbf{Transition.} À côté du SIEM, des sondes détectent au fil du réseau et
des hôtes~: le Chapitre~\ref{chap:idsips} traite des IDS/IPS.
\end{boxinfo}
